% ==============================================================================
% MÓDULO: CAPÍTULO V - ASPECTOS ADMINISTRATIVOS
% Archivo: subcapitulos/sec_05_aspectos_administrativos.tex
% ==============================================================================
\section{ASPECTOS ADMINISTRATIVOS}
\label{sec:cap5_administrativos}

\subsection{Asignación de recursos humanos y talento investigador}
\label{sec:5.1_recursos}
El desarrollo del proyecto de investigación estará sustentado por el siguiente equipo:
\begin{itemize}
    \item \textbf{Investigadora Principal:} Bach. Gresly Lucero Pariona Palomino, egresada de la Escuela Profesional de Enfermería de la Universidad Nacional de San Cristóbal de Huamanga (ORCID: 0009-0008-5421-9872). Responsable de la formulación teórica, trabajo de campo, realización de entrevistas, transcripción y redacción final de la tesis.
    \item \textbf{Asesor de Tesis:} Dr. Manglio Aguirre Andrade, Docente Principal y Catedrático de Proyecto de Investigación en Salud (EN 486) de la Facultad de Ciencias de la Salud -- UNSCH (ORCID: 0000-0001-8234-567X). Responsable de la orientación epistemológica, auditoría metodológica, supervisión ética y validación científica del estudio.
    \item \textbf{Personal de Apoyo Técnico y Lingüístico (02):} Asistentes de campo egresados de Enfermería con dominio fluido y acreditado del quechua chanka, responsables de la logística, apoyo en la preparación de ambientes de entrevista y co-verificación de la fidelidad de las transcripciones fonéticas.
\end{itemize}

\subsection{Presupuesto analítico y programación de gastos}
\label{sec:5.2_presupuesto}
El presupuesto general del proyecto ha sido proyectado a precios de mercado local en la ciudad de Ayacucho, desglosado en tres partidas presupuestarias esenciales:

\subsubsection{Pago de servicios de personal}
\label{sec:5.2.1_personal}
\begin{table}[H]
\centering
\small
\begin{tabularx}{\textwidth}{Xcccr}
\toprule
\textbf{Servicios Requeridos} & \textbf{U.M.} & \textbf{Cantidad} & \textbf{Costo Unit. (S/.)} & \textbf{Total (S/.)} \\
\midrule
Asesoría temática especializada y juicio de expertos & Mes & 04 & 500.00 & 2,000.00 \\
Personal de apoyo logístico y lingüístico quechua (02) & Mes & 02 & 500.00 & 1,000.00 \\
\midrule
\textbf{SUBTOTAL PARTIDA 1} & & & & \textbf{3,000.00} \\
\bottomrule
\end{tabularx}
\caption{Presupuesto analítico de servicios de personal.}
\label{tab:presupuesto_personal}
\end{table}

\subsubsection{Programación de otros servicios}
\label{sec:5.2.2_otros_servicios}
\begin{table}[H]
\centering
\small
\begin{tabularx}{\textwidth}{Xcccr}
\toprule
\textbf{Servicios Complementarios} & \textbf{U.M.} & \textbf{Cantidad} & \textbf{Costo Unit. (S/.)} & \textbf{Total (S/.)} \\
\midrule
Impresión de protocolos, consentimientos y guías & Hojas & 1000 & 0.20 & 200.00 \\
Conectividad y almacenamiento en la nube de audios & Horas & 200 & 1.00 & 200.00 \\
Licencia de software de análisis cualitativo (ATLAS.ti) & Unidad & 01 & 500.00 & 500.00 \\
Empastado formal de ejemplares según norma UNSCH & Unidad & 05 & 20.00 & 100.00 \\
Movilidad local y desplazamiento urbano a C.S. Belén & Días & 20 & 10.00 & 200.00 \\
\midrule
\textbf{SUBTOTAL PARTIDA 2} & & & & \textbf{1,200.00} \\
\bottomrule
\end{tabularx}
\caption{Presupuesto analítico de otros servicios.}
\label{tab:presupuesto_servicios}
\end{table}

\subsubsection{Adquisición de materiales y suministros de bioseguridad}
\label{sec:5.2.3_materiales}
\begin{table}[H]
\centering
\small
\begin{tabularx}{\textwidth}{Xcccr}
\toprule
\textbf{Materiales e Insumos} & \textbf{U.M.} & \textbf{Cantidad} & \textbf{Costo Unit. (S/.)} & \textbf{Total (S/.)} \\
\midrule
Material bibliográfico normativo y libros de metodología & Unidad & 02 & 50.00 & 100.00 \\
Papel Bond A4 80 gr. para trabajo de campo y borradores & Millar & 02 & 25.00 & 50.00 \\
Memoria USB de 32 GB con cifrado de seguridad para audios & Unidad & 02 & 40.00 & 80.00 \\
Archivadores, fólderes y separadores documentarios & Unidad & 10 & 1.00 & 10.00 \\
Útiles de escritorio (lapiceros, libretas, resaltadores) & Unidad & 10 & 2.00 & 20.00 \\
Mascarillas descartables KN95 para bioseguridad & Caja & 02 & 20.00 & 40.00 \\
Alcohol medicinal en gel al 70\% & Frasco & 02 & 10.00 & 20.00 \\
Fondo de reserva para imprevistos y contingencias operativas & Global & 01 & 80.00 & 80.00 \\
\midrule
\textbf{SUBTOTAL PARTIDA 3} & & & & \textbf{400.00} \\
\bottomrule
\end{tabularx}
\caption{Presupuesto analítico de materiales y suministros de bioseguridad.}
\label{tab:presupuesto_materiales}
\end{table}

\noindent\textbf{COSTO TOTAL ESTIMADO DEL PROYECTO: S/. 4,600.00 Soles.}

\subsection{Financiamiento}
\label{sec:5.3_financiamiento}
La ejecución integral del presente proyecto de investigación no demanda asignación presupuestaria a la Universidad Nacional de San Cristóbal de Huamanga ni al Ministerio de Salud, siendo asumido al 100\% con \textbf{recursos económicos propios de la investigadora principal}.

\subsection{Cronograma de actividades y dimensionamiento temporal (Técnica PERT)}
\label{sec:5.4_cronograma}
Para dimensionar con rigurosidad el tiempo operativo del proyecto se aplica la técnica de evaluación y revisión de programas (PERT), calculando el tiempo esperado ($t_e$) mediante la formulación tripartita de estimación temporal:
\begin{equation}
\label{eq:pert_tiempo}
t_e = \frac{O + 4m + P}{6}
\end{equation}
donde $O$ denota el tiempo optimista (4 meses), $m$ el tiempo más probable (6 meses) y $P$ el tiempo pesimista (10 meses), arrojando un horizonte medio de ejecución de 6.3 meses. 

En estricta observancia de los lineamientos del Dr. Manglio Aguirre Andrade, la planificación administrativa pondera la realidad burocrática del trámite universitario: tras la aprobación del proyecto en aula, el dictamen ante el Comité de Ética (CIEI-UNSCH) y la emisión de la resolución decanal toman aproximadamente de tres a seis meses. Por consiguiente, la fase de ejecución empírica de campo, procesamiento e informe final se programa operativamente entre los meses de agosto de 2026 y enero de 2027:

\begin{table}[H]
\centering
\small
\begin{tabularx}{\textwidth}{Xcccccc}
\toprule
\textbf{Fases y Actividades de Investigación} & \textbf{Ago 2026} & \textbf{Set 2026} & \textbf{Oct 2026} & \textbf{Nov 2026} & \textbf{Dic 2026} & \textbf{Ene 2027} \\
\midrule
Presentación y registro formal del proyecto & X & & & & & \\
Dictamen del Comité Institucional de Ética (CIEI) & & X & & & & \\
Validación de instrumentos (Juicio de expertos) & & X & & & & \\
Coordinaciones institucionales con C.S. Belén & & & X & & & \\
Trabajo de campo: Entrevistas a profundidad y notas & & & X & X & & \\
Transcripción verbatim y análisis con ATLAS.ti & & & & & X & \\
Validación de resultados con participantes & & & & & X & \\
Redacción y revisión del borrador de tesis con asesor & & & & & X & X \\
Sustentación formal de tesis de licenciatura & & & & & & X \\
\bottomrule
\end{tabularx}
\caption{Cronograma de actividades del proyecto de investigación (Técnica PERT).}
\label{tab:cronograma}
\end{table}

\subsection{Matriz de viabilidad institucional y gestión de riesgos operativos}
\label{sec:5.5_riesgos}
Para asegurar la continuidad ininterrumpida de la investigación frente a posibles contingencias en el Centro de Salud Belén, se formula la siguiente matriz preventiva:

\begin{table}[H]
\centering
\small
\begin{tabularx}{\textwidth}{p{3.2cm}Xp{5.5cm}}
\toprule
\textbf{Riesgo Operativo} & \textbf{Impacto Potencial} & \textbf{Estrategia de Mitigación Institucional} \\
\midrule
Desconfianza o temor del usuario a represalias & Negativa a participar o respuestas evasivas por miedo a perder su atención SIS. & Explicación personalizada del Consentimiento Informado en quechua; garantía absoluta de confidencialidad y código alfanumérico. \\
\addlinespace
Hacinamiento y falta de ambientes privados & Interrupciones y vulneración de la privacidad en el consultorio. & Coordinación anticipada con la Jefatura para uso de ambientes de biblioteca o telemedicina en horarios de baja afluencia. \\
\addlinespace
Rotación de personal médico en el establecimiento & Dificultades para la triangulación con profesionales tratantes. & Establecimiento de nexos institucionales con la Coordinación de Enfermería y médicos permanentes de la microred. \\
\bottomrule
\end{tabularx}
\caption{Matriz de contingencias y gestión de riesgos en campo.}
\label{tab:riesgos_operativos}
\end{table}

\newpage
